\section{A torus pencil with ordinary nodes}\label{sec:pencil}

We keep the notation of \Cref{prop:nonsplitting}.  In particular, curve
classes are identified with their Poincar\'e duals modulo torsion, and
$F$ is primitive, $F^2=0$, $[A]F=[C]F=0$, and $[B]F>0$.
The form $B$ will remain fixed throughout the construction.
There is already a unique curve image through every point.  To turn
these images into the fibers of a smooth map, we first remove critical
points of their sphere parametrizations using exact pair variations.
The normal equation will then supply foliated neighborhoods of the
embedded tori.  Finally, exact local changes near the remaining nodal
spheres will supply the corresponding charts at singular fibers.

\begin{theorem}[The pencil]\label{thm:pencil}
There are smooth one-forms $a_t,c_t$, $0\leq t\leq1$, with
$a_0=c_0=0$, such that
\[
 A_t=A+\dd a_t,\qquad C_t=C+\dd c_t
\]
and $(A_t,C_t,B)$ is a positive closed triple for every $t$.
For the structure $J$ determined by $(A_1,C_1)$ and tamed by $B$, the
irreducible curves of class $F$ are precisely the fibers of a smooth
proper map
\[
 p_X:X\longrightarrow S
\]
to a closed connected oriented surface.  Each regular fiber is an
embedded torus.  The singular fibers, if any, are finitely many immersed
rational curves, each with one transverse positive double point.

Every singular fiber has a neighborhood identified with a proper
holomorphic elliptic fibration $p_Y:Y\to\Delta$ with smooth total space,
a section, and one ordinary nodal central fiber.  In this identification
there are constants $x_i\in\R$, $y_i\in\R\setminus\{0\}$ such that
\begin{equation}\label{eq:pencil-model-pair}
 A_1=\operatorname{Re}\Sigma_i,\qquad
 C_1=x_i\operatorname{Re}\Sigma_i+y_i\operatorname{Im}\Sigma_i.
\end{equation}
Here $\Sigma_i$ is a nowhere-zero holomorphic two-form, and the model
can be chosen with punctured-disk presentation
\begin{equation}\label{eq:pencil-model-periods}
 z=a+\tau(s)b\pmod{\Z+\tau(s)\Z},\qquad
 \tau(s)=\frac{\log s}{2\pi i},\qquad
 \Sigma_i=h_i\,\dd s\wedge\dd z,
 \quad h_i\in\C\setminus\{0\}.
\end{equation}
Both $A_1$ and $C_1$ vanish on every fiber.
\end{theorem}

Only the qualitative existence of the holomorphic model in
\Cref{thm:nodal-model} is needed in this section.  Its construction is
local and independent of the pencil; none of the estimates involving
the collapsing parameter is used here.

\subsection{Curve compactness in the primitive class}

We use closed pseudoholomorphic curve theory for smooth almost complex
structures tamed by a symplectic form.  The particular inputs are
factorization of nonconstant maps into a simple map and a branched
cover \cite[Proposition~2.5.1]{McDuffSalamon2012}, positivity of
intersections in dimension four, and the adjunction formula with its
nonnegative weighted singularity defect
\cite[Theorem~2.3, Corollary~2.4, and Theorem~2.8]{Wendl2020}.
For compactness we use the genus-zero theory, including marked spheres
and their stable limits
\cite[Theorems~5.3.1 and~5.5.5]{McDuffSalamon2012}.
These results do not require regularity of the moduli space.  Our later
genericity argument must instead work within the allowed exact
variations of the pair, with $B$ fixed.

For every triple used below, \Cref{prop:nonsplitting} gives a curve of
class $F$ through every point, and every positive configuration of total
class $F$ consists of one simple image with multiplicity one.  Two
different such images are disjoint: an intersection would contribute
positively to their intersection number $F^2=0$.  A simple
parametrization $u:L\to X$ has
\begin{equation}\label{eq:pencil-adjunction}
 g(L)+\delta(u)=1.
\end{equation}
Consequently it is an embedded torus or a rational curve of defect one.
The latter may initially have a critical point.  The singularity
theorem for a simple curve implies that its critical and noninjective
points on the normalization form a finite set
\cite[Propositions~B.26 and~B.41 and Theorem~2.3]{Wendl2020}.
If it is immersed,
\Cref{eq:pencil-adjunction} says that it has exactly one double point,
and that its two branches meet transversely.

\begin{lemma}[Compactness without bubbling]\label{lem:pencil-compactness}
Let $(A_\nu,C_\nu,B_\nu)$ converge smoothly to a positive closed triple,
with all three cohomology classes equal to $[A],[C],[B]$, respectively.
Let $u_\nu:S^2\to X$ be simple curves of class $F$ for the structures
determined by $(A_\nu,C_\nu)$ and tamed by $B_\nu$.  After passage to a
subsequence and reparametrization, $u_\nu$ converges smoothly to a simple
sphere of class $F$.  The same statement holds with one marked point.
In particular, the absence of critical points on all such spheres is
open in a sufficiently high finite differentiability topology.
\end{lemma}

\begin{proof}
Let $B_\infty$ be the limiting third form.  It tames the limiting
structure with a positive lower bound on the unit tangent bundle of a
fixed metric.  Smooth convergence therefore makes this same fixed form
tame every sufficiently late structure.  Its area on each sphere is
$[B_\infty]F=[B]F$.  The fixed-tamer hypotheses and uniform energy
bound in the cited Gromov compactness theorem are consequently satisfied.
The nonconstant components in the stable
limit give a positive configuration of total class $F$ for the limiting
triple.  By \Cref{prop:nonsplitting}, there is exactly one simple image
and its total multiplicity is one.  In particular, the surviving map
is not a nontrivial cover.

The domain of a genus-zero stable limit is a tree of spheres.  With
only one nonconstant vertex, a nonempty tree of constant vertices would
have a constant leaf with at most one node and the one permitted mark.
It would have at most two special points and hence would be unstable.
There are therefore no constant components.  Gromov convergence is now
smooth convergence on the whole sphere after reparametrization, and
the marked points converge as well.  The same argument gives convergence
in the corresponding finite differentiability orders when the
backgrounds converge in a sufficiently high finite order.

If the limiting triple had no critical spheres but nearby triples had
such spheres, mark a critical point on each.  The conclusion just proved
would give a critical point on a limiting sphere.  This contradiction
proves openness.  This argument applies uniformly on any compact
parameter set of positive triples with the stipulated periods.
\end{proof}

\subsection{Transversality using exact variations of the pair}

The two closed pullback equations have two unavoidable integral
constraints.  We include the coefficient sphere in the universal
problem to account for them explicitly.

\begin{lemma}[Restricted first-jet transversality]\label{lem:pencil-jets}
An arbitrarily small smooth exact perturbation of $(A,C)$, with $B$
fixed and the triple positive, makes every simple sphere in class $F$
immersed.  The same assertion holds for a path of such pairs, relative
to immersed endpoints.

These perturbations may preserve any finite set of specified immersed
$F$-spheres that are common Lagrangians throughout the path.  At an
endpoint they may also leave fixed closed model subdisks already
foliated by $F$-curves.
\end{lemma}

\begin{proof}
We prove universal surjectivity in the allowed space of perturbations.
We reduce an annihilator to two constants arising from the pullback
periods and a distribution supported at the marked point.  The elliptic
symbol and first-jet tests remove the point-supported part; variations
of the coefficient sphere remove the two constants.
Fix the complex structure $j$ on $S^2$.  Write $v_0=(0,0,1)$ for the
coefficient direction of $B$.  For $v$ close to $v_0$, a convenient
closed anti-invariant pair for the structure $J_v$ is
\begin{equation}\label{eq:pencil-coefficient-pair}
 A_v=A-\frac{v_1}{v_3}B,\qquad
 C_v=C-\frac{v_2}{v_3}B.
\end{equation}
For any map $u$ in class $F$, the two pullback periods are
\[
 \left(\int_{S^2}u^*A_v,\int_{S^2}u^*C_v\right)
 =-[B]F\left(\frac{v_1}{v_3},\frac{v_2}{v_3}\right).
\]
Both pullbacks vanish when $u$ is $J_v$-holomorphic, so every
$F$-curve forces $v=v_0$.  Exact changes of $A,C$ leave the two
periods unchanged, whereas varying $v$ supplies both period directions
in the universal linearization.  Thus this auxiliary parameter removes
the period constraints without adding solutions in other coefficient
directions.  We retain it in the section
\begin{equation}\label{eq:pencil-universal-section}
 (u,z,v;A,C)\longmapsto
 \left(\overline\partial_{j,J_v}u,
       \partial_{j,J_v}u(z)\right).
\end{equation}
The second component is valued in the real rank-four bundle of complex
linear maps $T_zS^2\to T_{u(z)}X$.  On the zero set of the first
component, its vanishing is exactly $\dd u(z)=0$.

Use $W^{k,p}$ maps with $p>2$, $k\geq8$, and backgrounds of finite
class $C^r$, with $r\geq k+8$.  The first target is $W^{k-1,p}$.
Sobolev embedding gives $W^{k,p}\subset C^{k-1,\alpha}$ for some
$\alpha>0$, so first-jet evaluation at the moving mark is at least
$C^{k-2}$, hence $C^6$.  The background regularity gives the same
$C^6$ regularity for the composition maps in the first component.
Perturbing forms are $\dd\beta_A,\dd\beta_C$,
where the primitives belong to the $C^{r+1}$ closure of smooth
one-forms in the indicated support class.  This gives separable
Banach spaces.  We use the ordinary Banach bundle construction for
\Cref{eq:pencil-universal-section}; compare the first-jet framework in
\cite[Sections~2--3]{OhZhu2008}.  The surjectivity argument for our
restricted perturbations is given next.

Suppose $(\eta,\lambda)$ annihilates the derivative of
\Cref{eq:pencil-universal-section} at a simple critical sphere.
Here $\eta$ is a distribution dual to the Cauchy--Riemann target and
$\lambda$ is a covector on the jet target.  Testing map variations
supported away from $z$ gives $D_u^*\eta=0$ there.  Elliptic
regularity makes $\eta$ continuously differentiable away from $z$,
with the further finite regularity supplied by the coefficients.
The choice of $r$ gives more derivatives than are needed for all
distributional products and integrations by parts below.

\emph{Removal of tangential coefficients.}
Although the index uses a fixed $j$, we may test every infinitesimal
variation $\dot j$ of the domain complex structure.  On $S^2$ the map
$\xi\mapsto\mathcal L_\xi j$ is onto: its cokernel is
$H^{0,1}(TS^2)=0$, by the line-bundle surjectivity criterion since
$\deg TS^2=2>-2$ \cite[Theorem~3.23]{Wendl2014Lectures}.
Naturality under a domain diffeomorphism, with the
marked point moved by the inverse diffeomorphism, expresses the
$\dot j$ test as a combination of the already allowed map and mark
tests.  The direct variation of the jet with respect to $j$ vanishes
at $z$, because $\dd u(z)=0$.  Thus
\[
 \eta\left(\tfrac12J\dd u\,\dot j\right)=0
 \quad\text{for every }\dot j.
\]
At an immersion point these variations span the tangential
Cauchy--Riemann residuals.  Hence $\eta$ annihilates that tangential
subbundle away from $z$.

\emph{The two pullback coefficients.}
There are bundle maps $L_A,L_C$, with the regularity of the coefficients,
from Cauchy--Riemann residuals
to real two-forms on the domain such that, at a solution,
\begin{equation}\label{eq:pencil-pullback-linearization}
 \delta(u^*A_v)=L_A\,\delta(\overline\partial u),\qquad
 \delta(u^*C_v)=L_C\,\delta(\overline\partial u).
\end{equation}
These identities hold also for variations of the background pair and
of $v$.  To see that they extend without loss of regularity at a critical point, choose
$j\partial_s=\partial_t$ and put
$e=(u_s+Ju_t)/2$.  Anti-invariance gives the exact identity
\[
 A_v(u_s,u_t)=-2A_v(u_s,Je),
\]
and the same identity holds for $C_v$.  Differentiation at $e=0$
proves \Cref{eq:pencil-pullback-linearization}.  In particular the
maps $L_A,L_C$ contain $\dd u$, not its inverse, and extend continuously
across every critical point.  At an immersion point they vanish on
tangential residuals and identify the real two-dimensional normal
residual with the two pullback values: contraction with a nonzero
complex tangent vector by a nonzero complex volume form pairs the
complex normal line isomorphically with $\C$.

Consequently, on the injective immersion locus away from $z$, there
are unique continuously differentiable real functions $f_A,f_C$ such that
\begin{equation}\label{eq:pencil-annihilator-coefficients}
 \eta(e)=\int_{S^2}\bigl(f_A L_Ae+f_C L_Ce\bigr)
\end{equation}
for tests supported in that locus.  Let $U$ be a relatively compact
injective disk in the locus.  A neighborhood of its image can be
chosen to meet no other part of the curve.  Restrictions of smooth
ambient one-forms supported in this neighborhood are dense in
$\Omega^1_c(U)$ in the $C^1$ topology.  Indeed, the available regularity
of $u$ gives a compactly supported $C^2$ ambient extension of any such
source one-form in tubular coordinates.  Approximate this extension
in $C^1$ by smooth ambient one-forms, keeping their supports in the
same isolated neighborhood.  Their pullbacks converge in $C^1$.
Testing the variations $\dd\beta_A$ for these smooth primitives with
$u$ fixed and passing to the limit gives
\[
 \int_U f_A\,\dd\theta=0
 \quad\text{for every }\theta\in\Omega^1_c(U),
\]
since $f_A$ is continuous on the compact source support.
The jet contribution of this background variation is zero because
$\dd u(z)=0$.  Integration by parts proves $\dd f_A=0$ on $U$;
independent variations of $C$ prove $\dd f_C=0$.  The complement of
the finite exceptional set in $S^2$ is connected, so these functions
are global constants $k_A,k_C$ on that locus.

Define on the whole source the distribution with continuous coefficients
\[
 \eta_0(e)=k_A\int_{S^2}L_Ae+k_C\int_{S^2}L_Ce.
\]
Continuity of the $L$'s and elliptic regularity away from $z$ show
that $\eta-\eta_0$ is supported at $z$, including when there are other
critical or double points.  For every map variation $w$, closedness
of the two forms gives
\[
 \eta_0(D_uw)
 =k_A\int_{S^2}\dd\bigl(u^*\iota_wA\bigr)
  +k_C\int_{S^2}\dd\bigl(u^*\iota_wC\bigr)=0.
\]

\emph{Elimination of the point-supported remainder.}
Set $R=\eta-\eta_0$.  The annihilator equation for map variations is
\begin{equation}\label{eq:pencil-point-distribution}
 D_u^*R=-\mathcal J_z^*\lambda,
\end{equation}
where $\mathcal J_z$ is the linearization of the complex-linear first
jet, with the mark fixed.  The right side has distributional order at
most one and is supported at $z$.  A point-supported distribution is
a finite sum of derivatives of the Dirac delta.  If its largest
nonzero order were $m\geq1$, its homogeneous coefficient polynomial
$R_m(\xi)$ would satisfy
\[
 \sigma(D_u)^*(\xi)R_m(\xi)=0
 \quad(\xi\in T_z^*S^2).
\]
There are enough coefficient derivatives for this comparison: a
point-supported functional on $W^{k-1,p}$ in real dimension two has
order at most $k-2$ when $p>2$, while $r\geq k+8$.
The first-order Cauchy--Riemann symbol is invertible for every real
$\xi\ne0$, so this polynomial must vanish identically, a
contradiction.  Hence $R=q\delta_z$ for a single covector $q$.
Test \Cref{eq:pencil-point-distribution} on variations with $w(z)=0$.
Since $\dd u(z)=0$, the values of $D_uw$ and $\mathcal J_zw$ at
$z$ are respectively the complex-antilinear and complex-linear parts
of $\nabla w(z)$.  These parts can be specified independently by a
smooth variation supported near $z$.  It follows that $q=0$ and
$\lambda=0$.

Finally, variation of $v$ in \Cref{eq:pencil-coefficient-pair} gives
\begin{equation}\label{eq:pencil-period-cokernel}
 \delta\left(\int_F A_v,\int_F C_v\right)
 =-[B]F\,(\dot v_1,\dot v_2).
\end{equation}
This map is onto $\R^2$.  Applying the remaining annihilator
$\eta_0$ to these tests yields $k_A=k_C=0$.  Thus the universal
derivative is onto.  Its range is closed: the map-and-jet operator
with background fixed is Fredholm and has finite-codimensional
closed range, and adding parameter directions preserves this
property.

\emph{Index and allowed genericity.}
The parametrized sphere operator has real index $4$, since
$c_1(u^*TX)=0$ \cite[Theorem~3.22]{Wendl2014Lectures}.
Adding $v$ and the mark and imposing the first-jet
constraint gives
\begin{equation}\label{eq:pencil-critical-index}
 4+2+2-4=4.
\end{equation}
For a path, allowing its parameter adds one and gives index $5$.
The projection from the universal zero set to the space of perturbing
primitives is therefore Fredholm of index $4$, or $5$ for paths.
The projection is $C^6$, and $6>5$, so Sard--Smale applies with the
chosen regularities \cite[Theorem~1.3]{Smale1965}.  For a regular
parameter its fiber would be a manifold of that dimension.  But the
six-dimensional group $\operatorname{PSL}(2,\C)$ acts freely on
simple parametrized spheres, with the marked point transformed
contragrediently, and preserves the fiber.  Freeness follows because
a reparametrization fixing a simple map fixes its injective locus and
hence is the identity.  A manifold of dimension $4$ or $5$ cannot
contain this six-dimensional orbit.  The fiber is empty.

For a path, the local perturbations just used may be multiplied by
arbitrary bumps in the path parameter, supported away from its
endpoints.  If finitely many immersed $F$-spheres are retained, a
critical simple $F$-sphere cannot have one of those images.  It is
therefore disjoint from all of them by positivity of intersections,
so every injective disk needed in the proof still admits the allowed
local variations.  Use the closure of smooth primitives supported
off the retained images; this preserves them exactly without requiring
a fixed forbidden tubular neighborhood.  At a fixed model subdisk,
every point already lies on a model $F$-curve.  A competing critical
sphere is disjoint from the whole subdisk for the same reason.  The
endpoint assertion follows with perturbations supported in its
complement.

There are countably many map-space charts and homotopy components,
so the regular parameter choices can be made simultaneously for all
of them.  The conclusions are open by \Cref{lem:pencil-compactness}.
Approximation in the defining closures by smooth primitives therefore
gives smooth perturbations with the same conclusion and the same
support restrictions.  For paths with immersed endpoints, openness
allows fixed small parameter collars at both ends.  Smallness ensures
positivity with $B$ fixed; the straight segment realizes the initial
small perturbation by an exact positive-pair path.
\end{proof}

\Needspace{13\baselineskip}
\subsection{Deforming immersed curves}

The identification of normal deformations with one-forms follows the
approach of McLean's deformation theory for compact special Lagrangians
\cite[Theorem~3.6 and Corollary~3.9]{McLean1998}.
Here the definite pair leads to a variable-coefficient operator, which
we solve directly.

\begin{lemma}[The normal operator]\label{lem:pencil-normal}
Let $u:L\to X$ be a closed connected common Lagrangian immersion for a definite
closed pair $(A,C)$, with the complex orientation chosen by $B$.
Nearby common Lagrangian immersions represented as normal graphs over
$u$ form a smooth manifold of real dimension $b_1(L)$.
Its linearized kernel is naturally isomorphic to $H^1(L;\R)$.
If $u$ is simple, this also describes the local unparametrized moduli.
These descriptions hold with smoothly varying pairs whose periods on
$[u(L)]$ remain zero.

An immersed $F$-sphere is isolated and continues uniquely under small
pair variations.  Near an embedded $F$-torus the nearby $F$-curves
form a smooth foliation of an open neighborhood.
\end{lemma}

\begin{proof}
Represent nearby immersions by normal graphs over $u$, removing the
infinitesimal reparametrizations.  If $u$ is simple, its stabilizer under
reparametrization is trivial: an automorphism fixing the map fixes its
injective locus and hence is the identity.  The normal-graph slice then
also parametrizes the local unparametrized moduli.  For a multiply
covered immersion we make the smoothness assertion only for this slice.
The symplectic form $A$ identifies the real normal bundle with $T^*L$
by $w\mapsto\alpha=u^*\iota_wA$.  Write along the source
\begin{equation}\label{eq:pencil-volume-normalization}
 H=\frac{C-xA}{y},\qquad \Omega=A+iH,\qquad
 x=\frac{A\wedge C}{A^2},\qquad
 y^2=\frac{C^2}{A^2}-x^2,
\end{equation}
where the nonzero sign of $y$ is the one giving $J$ tamed by $B$.
Let $*_L$ be the star on one-forms for the induced conformal structure
and orientation, with $*_L\alpha=-\alpha\circ j$.
The complex volume identity gives
\[
 u^*\iota_wC=x\alpha+y*_L\alpha.
\]
Cartan's formula now gives the normal derivative
\begin{equation}\label{eq:pencil-normal-operator}
 \mathcal D\alpha
   =\left(\dd\alpha,\dd(x\alpha+y*_L\alpha)\right).
\end{equation}
The nonlinear pullbacks and their derivatives take values in pairs
of exact two-forms: their integrals are the two fixed zero periods.

For completeness we solve \Cref{eq:pencil-normal-operator}.  Given
exact two-forms $\zeta_1,\zeta_2$, choose $\alpha_0$ with
$\dd\alpha_0=\zeta_1$.  Adding $\dd f$ changes only the second
equation, by
\begin{equation}\label{eq:pencil-scalar-operator}
 \mathcal P f
 =\dd(x\dd f+y*_L\dd f)
 =\dd x\wedge\dd f+\dd(y*_L\dd f).
\end{equation}
After division by a positive area form its principal part is the
nonzero signed multiple $y\Delta f$ of the scalar Laplacian.
Its sign is constant on connected $L$, it has no zeroth-order term,
and the strong maximum principle shows that its kernel consists exactly
of constants \cite[Chapter~1, Theorem~1.6]{KiselevRoquejoffreRyzhik2012}.
Here one first adjusts the overall sign of the operator and then applies
the principle at a maximum on the connected closed surface.
It has index zero, by homotopy through scalar elliptic operators to the
signed Laplacian, with the sign of $y$ held fixed; the Fredholm index is
constant along this homotopy
\cite[Section~7, Propositions~7.3--7.4]{TaylorFunctionalNotes}.
Its image has integral zero by Stokes' theorem.  Its cokernel has
dimension one, so its range is exactly the integral-zero two-forms.
We can therefore solve
\[
 \mathcal P f=\zeta_2-\dd(x\alpha_0+y*_L\alpha_0).
\]
This proves surjectivity onto the asserted target, with bounded
elliptic right inverses in the usual Sobolev or H\"older spaces.

If $\mathcal D\alpha=0$, then $\alpha$ is closed.  Every class in
$H^1(L;\R)$ has a unique representative of this kind: start with any
closed representative and solve \Cref{eq:pencil-scalar-operator} for
an exact correction.  Uniqueness follows from the constant kernel of
$\mathcal P$.  Thus
\begin{equation}\label{eq:pencil-kernel-cohomology}
 \ker\mathcal D\longrightarrow H^1(L;\R),\qquad
 \alpha\longmapsto[\alpha]
\end{equation}
is an isomorphism.  The implicit function theorem applied to the
pullback equations, with exact two-forms as target, proves the first
assertion and its parameter version.  For $L=S^2$ the derivative is
an isomorphism, giving isolation and unique continuation.

For an embedded torus, the kernel has dimension two.  In complex
normal notation its elements solve the real-linear normal
Cauchy--Riemann equation.  Indeed, the normal projection of a
linearized map equation is unchanged by variation of the source
complex structure, and the two maps in
\Cref{eq:pencil-pullback-linearization} identify that normal equation
with \Cref{eq:pencil-normal-operator}.  The complex normal line has
degree $F^2=0$.  The similarity principle says that a nonzero solution
has isolated positive zeros, whose total number with multiplicity is
the degree; in the closed case this is the zero formula of
\cite[Section~2.2, Equation~(2.7)]{Wendl2010}.  Therefore every nonzero
kernel section is nowhere vanishing.  Evaluation at each point of the
torus is consequently an isomorphism from the two-dimensional kernel
to its real normal plane.

Let $u_q$, $q\in\R^2$ small, be the family supplied by the implicit
function theorem.  The derivative of $(x,q)\mapsto u_q(x)$ is an
isomorphism at every $(x,0)$.  Compactness of the torus gives a common
small parameter disk on which this is a local diffeomorphism.  The
individual maps remain embeddings; different images cannot meet by
$F^2=0$ and positivity.  The immersion slice distinguishes the images,
so evaluation is injective after shrinking the disk.  It is thus a
diffeomorphism onto an open tube, foliated by the $u_q(L)$.
\end{proof}

\begin{corollary}[The number of nodes along a path]
\label{cor:pencil-node-count}
If all $F$-spheres for a triple are immersed, there are finitely many
of them.  Along a compact smooth path of triples with the original
periods and with all $F$-spheres immersed, their number is constant.
\end{corollary}

\begin{proof}
The unparametrized sphere space is compact by
\Cref{lem:pencil-compactness} and discrete by
\Cref{lem:pencil-normal}, so it is finite.  Over a path, the total
space of such spheres is compact and, by the parameter version of
\Cref{lem:pencil-normal}, its projection to the path interval is a
local diffeomorphism, also in relative half-interval charts at the
endpoints.  Hence it is a finite covering of the interval.  Equivalently,
each isolated sphere has a unique local continuation, and compactness
prevents any additional sphere from entering those local descriptions.
The number is constant.
\end{proof}

\subsection{Exact insertion of holomorphic nodal neighborhoods}

\begin{lemma}[Insertion while retaining the sphere]
\label{lem:pencil-insertion}
Let $K$ be the image of an immersed $F$-sphere with one transverse
positive double point.  Inside an arbitrarily small neighborhood of
$K$, an exact path of the pair, with $B$ fixed, makes the pair equal
near $K$ to \Cref{eq:pencil-model-pair} for a model satisfying
\Cref{eq:pencil-model-periods}.  The triple stays positive and the
same immersed sphere is a common Lagrangian throughout.  The
construction may be performed in disjoint neighborhoods of finitely
many such spheres.
\end{lemma}

\begin{proof}
Denote the node by $p$ and its two preimages on the normalization by
$p_+,p_-$.  We first identify the model with a neighborhood of $K$
so that straight interpolation of the two pairs stays positive.
We then localize this interpolation by two exact changes.  At the node,
equality of the forms as tensors makes the radial cutoff correction small.  Along
the rest of the sphere, the difference of the pairs need not be small;
a logarithmic transverse cutoff makes its cutoff error small while
retaining the positive interpolation.

Use $x,y,\Omega$ from
\Cref{eq:pencil-volume-normalization} on a neighborhood of $K$.
Choose the qualitative model of \Cref{thm:nodal-model}; write its
normalization as $u_Y:S^2\to Y$ with node preimages $q_+,q_-$.
Uniformization supplies a biholomorphism $f$ of the normalized
spheres carrying $p_\pm$ to $q_\pm$.

\emph{Matching the derivative and the complex volume.}
At $p$, the two complex tangent lines of the branches are transverse.
Their prescribed tangent maps, obtained from $f$, define one complex
linear isomorphism
\[
 T:T_pX\longrightarrow T_{p_Y}Y.
\]
Choose the nonzero complex scale $h$ of the model volume so that
$T^*\Sigma_Y=\Omega(p)$.  Along each smooth point of the
normalization, the tangent map and equality of complex volumes
determine the map on the normal quotient.  More explicitly, a
nonvanishing complex volume identifies the complex normal line with
the dual of the tangent line.  The prescribed tangent map therefore
gives the required normal-line isomorphism globally.  The possible
lifts to a map of the two ambient complex plane bundles form an
affine bundle with fiber the complex-linear maps from the normal
line to the target tangent line.  A smooth section exists by
partitions of unity.  Its values at both preimages of the node can
be chosen to equal the single map $T$.

These prescribed first derivatives need only be realized to arbitrary
accuracy along $K$, with exact realization at the node.  Here is a
construction that addresses the crossing.  First straighten the two
transverse branches by smooth ambient charts in the source and target.
Writing the restrictions of $f$ in these charts as $f_+(u)$ and
$f_-(v)$, the product map
$\Phi_0(u,v)=(f_+(u),f_-(v))$ is a genuine local diffeomorphism
carrying the branches by $f$ and having derivative $T$ at $p$.
Its derivative along either branch
and the prescribed complex-linear derivative have the same tangent
restriction and agree at $p$.  They are consequently arbitrarily
close on sufficiently short branch collars.  On those collars
interpolate their normal derivative data, keeping $\Phi_0$ near
$p$ and the prescribed data away from the collars.  The interpolated
derivatives remain invertible when the collars are sufficiently
short.

Away from the crossing the normal tubular construction realizes these
jets: in normal coordinates send the zero section by $f$ and use the
prescribed normal derivative on the normal coordinate.  This can be
patched to $\Phi_0$ where the jets already agree; a partition in
normal coordinates preserves the prescribed first jets.  Shrinking
the tubes gives a diffeomorphism $\Phi$ of neighborhoods of the nodal
images.  To check injectivity, any hypothetical pair of coincident
images in arbitrarily shrinking neighborhoods would limit to two
points of $K$ with the same image.  The map on the nodal image is a
homeomorphism, and the map is a local diffeomorphism at every point,
including the crossing.  These facts rule out such a pair.
We have arranged that
\begin{equation}\label{eq:pencil-volume-matching}
 \Phi^*\Sigma_Y=\Omega\quad\text{at }p,
 \qquad
 \Phi^*\Sigma_Y\text{ is arbitrarily close to }\Omega
 \quad\text{along }K.
\end{equation}
Both volumes vanish when pulled back to the normalization.

Put $x_p=x(p)$, $y_p=y(p)$ and set on this neighborhood
\[
 A_*=\operatorname{Re}\Phi^*\Sigma_Y,\qquad
 C_*=x_p\operatorname{Re}\Phi^*\Sigma_Y
       +y_p\operatorname{Im}\Phi^*\Sigma_Y.
\]
These forms are closed.  If the volume matching along $K$ were exact,
their pointwise relation to the old pair would be
\begin{equation}\label{eq:pencil-triangular-matching}
 A_*=A,\qquad
 C_*=\left(x_p-\frac{y_px}{y}\right)A+\frac{y_p}{y}C.
\end{equation}
The ratio $y_p/y$ is positive.  Thus every straight interpolation of
the two pairs in \Cref{eq:pencil-triangular-matching} has the same
positive anti-invariant plane and remains in the same component
relative to $B$.  In particular its triple Gram matrix has a uniform
positive lower bound along $K$, with the interpolation parameter in
$[0,1]$.  Choose the approximation in
\Cref{eq:pencil-volume-matching} small compared with this bound.
The actual straight interpolation between $(A,C)$ and $(A_*,C_*)$
is then positive along $K$, and on a smaller neighborhood by
compactness and continuity.
The biholomorphism $f$ uses the complex orientation selected by $B$,
so $B$ is positive on the transported model's complex tangent line at
each smooth point of $K$.  Together with positivity of the triple,
this selects the model complex structure as the sign tamed by $B$.
That sign persists on a smaller connected model neighborhood,
independently of whether $y_p$ is positive or negative.

\emph{The first exact correction, at the crossing.}
For either closed difference $\xi=A_*-A$ or $\xi=C_*-C$, one has
$\xi(p)=0$ as a tensor, and the pullback of $\xi$ to each branch
vanishes.  In a crossing chart with $p=0$, the radial homotopy
primitive is
\begin{equation}\label{eq:pencil-radial-primitive}
 \beta_z(v)=\int_0^1 t\,\xi_{tz}(z,v)\,\dd t,
 \qquad \dd\beta=\xi.
\end{equation}
It satisfies $|\beta(z)|\leq C|z|^2$.  Since radial contraction
preserves each straightened branch, $\beta$ pulls back to zero on
both branches.  Choose $\chi_r$ equal to one on the ball of radius
$r$, supported in the ball of radius $2r$, and satisfying
$|\dd\chi_r|\leq C/r$.  The exact change
$\dd(\chi_r\beta)$ has norm $O(r)$: the two terms
$\chi_r\xi$ and $\dd\chi_r\wedge\beta$ both have that bound.
For sufficiently small $r$ its scalar multiples preserve positivity
with $B$ fixed.  They preserve the common Lagrangian sphere because
the primitive restricts to zero there.  Do this for both differences.
The resulting pair $(A',C')$ equals $(A_*,C_*)$ on the ball of
radius $r$ and differs from $(A,C)$ by $O(r)$ near that ball.

\emph{The exact correction along the rest of the sphere.}
The remaining differences $\xi'=A_*-A'$ and $\xi'=C_*-C'$ are
closed, vanish on a neighborhood of the crossing, and pull back to
zero on the source.  Outside a smaller crossing ball, take normal
tubes of the embedded part of the source.  Normal contraction
$h_t(s,n)=(s,tn)$ gives a relative primitive $\beta'$ satisfying
\begin{equation}\label{eq:pencil-tubular-primitive}
 \dd\beta'=\xi',\qquad
 |\beta'|\leq C d,
 \qquad d=|n|.
\end{equation}
This is the homotopy identity
$\dd\beta'=\xi'-\pi^*u^*\xi'$.  Choose the tubes in the inner
collars so that the contraction stays where $\xi'=0$.  Then
$\beta'=0$ on those collars and extends by zero into the crossing
ball.  The two branch tubes are disjoint outside that ball after
shrinking.  This gives a single smooth relative primitive on the
required tubular region, without a compatibility condition left at
the node.

Choose a transverse cutoff equal to one for $d\leq r_{\rm in}$ and
zero for $d\geq r_{\rm out}$, where the outer tube is already small.
Using a fixed smooth function of
$\log(d/r_{\rm in})/\log(r_{\rm out}/r_{\rm in})$ gives
\begin{equation}\label{eq:pencil-log-cutoff}
 d\,|\dd\chi|\leq
 \frac{C}{\log(r_{\rm out}/r_{\rm in})}.
\end{equation}
The inner radius can be made sufficiently small that the right
side is any prescribed $\eta>0$.  For $0\leq t\leq1$, the exact
change is
\[
 A'+t\dd(\chi\beta'_A)
  =A'+t\chi(A_*-A')+t\dd\chi\wedge\beta'_A,
\]
and likewise for $C'$.  The last terms have norm at most $C\eta$
by \Cref{eq:pencil-tubular-primitive,eq:pencil-log-cutoff}.
The first terms are pointwise straight interpolation with parameter
$t\chi\in[0,1]$.  Taking $r$ small in the first correction preserves
the positive margin already established along $K$.  Shrinking the
outer tubes preserves that margin throughout them.  Finally take
$\eta$ small compared with the margin.  These choices prove
positivity along the entire exact path.

The final pair agrees with the model on a neighborhood of the whole
nodal image, and its primitives have compact support in the original
chosen neighborhood.  A sufficiently small complete model subdisk
lies in this region: properness of $p_Y$ and compactness of its
central fiber imply that $p_Y^{-1}(\Delta_\rho)$ lies in any chosen
neighborhood of that fiber when $\rho$ is sufficiently small.
Every modification restricts to zero on the normalization, so the
same sphere is retained.  Disjoint supports prove the last assertion.
\end{proof}

\subsection{Retaining the nodes and forming the base}

\begin{proof}[Proof of \Cref{thm:pencil}]
First apply \Cref{lem:pencil-jets} to make a small exact perturbation
with $B$ fixed so that all $F$-spheres are immersed.  Connect it to
the starting pair by a small straight segment of positive pairs.
There is no assertion about immersion along this initial segment.
By \Cref{cor:pencil-node-count} the perturbed pair has finitely many
nodal spheres $K_1,\ldots,K_N$, all disjoint.  If $N=0$, the torus
argument below already applies.

For $N>0$, use \Cref{lem:pencil-insertion} at all the $K_i$, keeping
$B$ fixed and retaining every $K_i$ along the resulting path.
Choose smaller closed model subdisks at its endpoint.  Make the
endpoint generic by an arbitrarily small exact perturbation outside
those subdisks.  The relative assertion of \Cref{lem:pencil-jets}
applies because any competing $F$-curve is disjoint from every
model fiber.  The endpoint now has no critical spheres, and the
smaller model neighborhoods remain unchanged.  Extend this endpoint
perturbation a short distance into the path using a parameter cutoff;
all retained $K_i$ are still common Lagrangians.

The two ends of this latter path have only immersed spheres.  Apply
the path version of \Cref{lem:pencil-jets}, keeping its endpoints
and every $K_i$ fixed, to remove all critical spheres from its
interior.  The resulting smooth path stays positive with $B$ fixed.
By \Cref{cor:pencil-node-count}, its endpoints have the same number
of nodal spheres.  All the original $N$ spheres were retained, so
they exhaust the endpoint spheres.  This argument is needed to
exclude additional immersed spheres outside the inserted model
regions.

We now work at the endpoint.  By \Cref{prop:nonsplitting}, every
point of $X$ lies on a unique simple $F$-curve image.  Each is an
embedded torus or one of the $K_i$.  Define $S$ as the set of these
images, with quotient map $p_X$.  For a torus,
\Cref{lem:pencil-normal} supplies an open tube diffeomorphic to
$T^2\times\Delta$.  It is saturated: any other $F$-curve meeting
a torus in this tube must equal it by intersection positivity.
For a nodal sphere choose an open model subdisk.  Its fibers also
represent $F$, since they are homologous to its central fiber.
The same argument makes this open subset saturated.  Its quotient
is its base disk, since the proper holomorphic model projection is
open and has connected fibers.  Thus the torus tubes and the model
subdisks give disk neighborhoods in $S$ and cover it.

These saturated neighborhoods can be chosen arbitrarily small
around an entire fiber.  For tori this follows by shrinking the
parameter disk in the tube; for nodes it follows from properness of
the model projection.  Two distinct fibers are disjoint compact
subsets of $X$, so they admit disjoint ambient neighborhoods.
Choose saturated neighborhoods inside them.  Their images are
disjoint open neighborhoods of the two points of $S$.  The quotient
is therefore Hausdorff.

The disk charts have smooth transition maps.  To verify this even
at a singular value, choose a smooth point of the corresponding
fiber.  Near that point both local quotient projections are smooth
submersions with the same fibers.  A local section of either
submersion expresses the other base coordinate as a smooth function
of it; interchanging the two gives its smooth inverse.  The charts
therefore make $S$ a smooth surface without boundary and $p_X$ a
smooth map.  They also make $p_X$ open.  Since $X$ is second
countable, images of a countable open basis under this open
surjection form a countable basis for $S$.  Compactness and
connectedness of $X$ give compactness and connectedness of $S$.

On regular fibers use their complex orientation and the complex
orientation of the normal plane to orient the base.  Evaluation
along a connected torus preserves this choice, and the holomorphic
model gives the same orientation on its regular part.  Hence the
orientation extends across each singular value.  The map is proper
because its domain is compact and its base is Hausdorff.  Its
singularities are precisely the ordinary holomorphic nodes of the
models (locally $s=z_1z_2$ after complex coordinate changes).
The pair vanishes on every fiber by construction, and the prescribed
model identities hold on the unchanged smaller neighborhoods.

All paths used above consist of exact changes to $A,C$ and keep
$B$ fixed.  Their finitely many joins can be made smooth by
reparametrizations constant to infinite order at the join times.
Concatenating them proves the theorem.
\end{proof}
